% 第16批：独立推导、自查；依赖源图的题标明判读前提。
\subsection{自测题5：统计物理、热力学计算与证明}

\sourcetest{5.3.1}{117}{125}
\problemtext 体积1m$^3$ 中有 $N_1=1.0\times10^{25}$ 个氧分子和 $N_2=4.0\times10^{25}$ 个氮分子，混合气体压强 $2.76\times10^5$Pa；求每分子平均平动动能及温度。
\workings 总粒子数 $N=5.0\times10^{25}$，由 $pV=Nk_BT$ 得 $k_BT=pV/N=5.52\times10^{-21}$J。平均平动能 $\boxed{\bar\varepsilon_{\rm tr}=3k_BT/2=8.28\times10^{-21}\ {\rm J}}$；取 $k_B=1.380649\times10^{-23}$J/K，$\boxed{T\simeq399.8\ {\rm K}}$。
\crosscheck 混合气体平衡态两种分子的每粒子平均平动能相同，与分子质量无关。

\sourcetest{5.3.2}{117}{125}
\problemtext 氢核（质子）温度达到 $10^8$K，以非相对论理想气体近似求方均根速率与平均平动能（eV）。
\workings $\boxed{v_{\rm rms}=\sqrt{3k_BT/m_p}\simeq1.57\times10^6\ {\rm m/s}}$，取 $m_p=1.673\times10^{-27}$kg。$\bar\varepsilon=3k_BT/2=2.071\times10^{-15}$J；以 $1\,{\rm eV}=1.602176634\times10^{-19}$J，$\boxed{\bar\varepsilon\simeq1.29\times10^4\,{\rm eV}=12.9\,{\rm keV}}$。
\crosscheck $v_{\rm rms}\ll c$，非相对论温度速度估计自洽。

\sourcetest{5.3.3}{117}{125}
\problemtext 15°C、压强0.76mHg时氩气和氖气平均自由程分别为 $6.7\times10^{-8}$m、$13.2\times10^{-8}$m；求有效碰撞直径比 $d_{\rm Ne}/d_{\rm Ar}$；氩在20°C、0.15mHg的平均自由程。
\workings 硬球气体 $\lambda=k_BT/(\sqrt2\pi d^2p)$，同态下 $\boxed{d_{\rm Ne}/d_{\rm Ar}=\sqrt{\lambda_{\rm Ar}/\lambda_{\rm Ne}}=\sqrt{6.7/13.2}\simeq0.712}$。同种氩气改变状态时 $\lambda'_{{\rm Ar}}=\lambda_{\rm Ar}(T_2/T_1)(p_1/p_2)$；$T_1=288.15$K，$T_2=293.15$K，故 $\boxed{\lambda'_{{\rm Ar}}\simeq6.7\times10^{-8}(293.15/288.15)(0.76/0.15)\simeq3.45\times10^{-7}\ {\rm m}}$。
\crosscheck 两压强单位均为米汞柱，可直接取比；有效直径是简化硬球碰撞参数。

\sourcetest{5.3.4}{117--118}{NaN}
\problemtext 2mol氢气初态温度27°C、体积20L，先等压膨胀使体积加倍，再可逆绝热膨胀至温度回到初温；求 $pV$ 过程、吸热、内能变化与总功。
\workings 以刚性双原子 $\gamma=7/5$ 计算。$T_1\simeq300$K、$V_1=20$L；等压末态 $V_2=40$L、$T_2=600$K；随后绝热末体积 $V_3=V_2(T_2/T_3)^{1/(\gamma-1)}=40\cdot2^{5/2}\simeq226.3$L。等压吸热 $Q_{12}=n(7R/2)(T_2-T_1)=2(7R/2)(300)\simeq\boxed{1.75\times10^4\ {\rm J}}$，绝热段 $Q_{23}=0$，总 $\Delta U=nC_V(T_3-T_1)=\boxed{0}$；故全过程 $W=Q-\Delta U=\boxed{1.75\times10^4\ {\rm J}}$。
\crosscheck 这是开放路径而非热机循环，$T_3=T_1$ 不代表 $V_3=V_1$。

\sourcetest{5.3.5}{118}{126}
\problemtext 氮气初始1atm、300K，经可逆绝热压缩至32atm；求末温及末态粒子数密度。
\workings 取 $\gamma=7/5$；$T_2=T_1(p_2/p_1)^{(\gamma-1)/\gamma}=300\cdot32^{2/7}\simeq\boxed{807.6\ {\rm K}}$。数密度 $n_N=p_2/(k_BT_2)$，用 $1$atm$=101325$Pa，$\boxed{n_N\simeq(32\cdot101325)/(1.380649\times10^{-23}\cdot807.6)\simeq2.91\times10^{26}\ {\rm m^{-3}}}$。
\crosscheck 这里 $n_N$ 是每立方米分子数，并非摩尔数；绝热指数假设分子振动自由度未激发。

\sourcetest{5.4.1}{118}{126}
\problemtext 证明两恒温热源间可逆卡诺热机效率 $\eta=1-T_c/T_h$。
\workings 卡诺循环由温度 $T_h,T_c$ 的两段可逆等温与两段可逆绝热组成。可逆绝热段 $\Delta S=0$，故两等温段气体熵变大小相等：$Q_h/T_h=-Q_c/T_c$，其中 $Q_h>0$、$Q_c<0$。整循环 $\Delta U=0$，输出功 $W=Q_h+Q_c$，所以 $\boxed{\eta=W/Q_h=1+Q_c/Q_h=1-T_c/T_h}$。
\crosscheck 只对两个恒温热源间可逆卡诺机成立；不可逆实际热机效率通常更低。
